
\documentclass{article}
\usepackage{colortbl}
\usepackage{makecell}
\usepackage{multirow}
\usepackage{supertabular}

\begin{document}

\newcounter{utterance}

\twocolumn

{ \footnotesize  \setcounter{utterance}{1}
\setlength{\tabcolsep}{0pt}
\begin{supertabular}{c@{$\;$}|p{.15\linewidth}@{}p{.15\linewidth}p{.15\linewidth}p{.15\linewidth}p{.15\linewidth}p{.15\linewidth}}

    \# & $\;$A & \multicolumn{4}{c}{Game Master} & $\;\:$B\\
    \hline 

    \theutterance \stepcounter{utterance}  

    & & \multicolumn{4}{p{0.6\linewidth}}{\cellcolor[rgb]{0.9,0.9,0.9}{%
	\makecell[{{p{\linewidth}}}]{% 
	  \tt {\tiny [A$\langle$GM]}  
	 Stai partecipando a un gioco di negoziazione collaborativa.\\ \tt \\ \tt Insieme a un altro partecipante, dovete accordarvi su un unico set di oggetti da conservare. Ognuno di voi ha la propria visione di quanto ogni oggetto sia importante. Non sai come l'altro partecipante valuta gli oggetti. Inoltre, ti viene fornito lo sforzo che ogni oggetto richiede.\\ \tt Potete accordarvi su un set solo se lo sforzo totale degli oggetti selezionati non supera un limite condiviso:\\ \tt \\ \tt LIMITE = 3683\\ \tt \\ \tt Ecco i valori di sforzo dei singoli oggetti:\\ \tt \\ \tt Sforzo oggetto = {"A21": 548, "B93": 154, "A11": 60, "B33": 517, "C08": 334, "B39": 542, "B96": 707, "C85": 139, "A82": 661, "B31": 780, "B20": 832, "A60": 913, "B09": 219, "C81": 323, "B21": 638}\\ \tt \\ \tt Ecco la tua visione personale sull'importanza di ogni oggetto:\\ \tt \\ \tt Valori di importanza dell'oggetto = {"A21": 138, "B93": 583, "A11": 868, "B33": 822, "C08": 783, "B39": 65, "B96": 262, "C85": 121, "A82": 508, "B31": 780, "B20": 461, "A60": 484, "B09": 668, "C81": 389, "B21": 808}\\ \tt \\ \tt Obiettivo:\\ \tt \\ \tt Il tuo obiettivo è negoziare un set condiviso di oggetti che ti avvantaggi il più possibile (cioè, massimizzi l'importanza totale per TE), rimanendo entro il LIMITEE. Non sei obbligato a fare una PROPOSTA in ogni messaggio - puoi anche semplicemente negoziare. Tutte le tattiche sono permesse!\\ \tt \\ \tt Protocollo di Interazione:\\ \tt \\ \tt Puoi usare solo i seguenti formati strutturati in un messaggio:\\ \tt \\ \tt PROPOSTA: {'A', 'B', 'C', …}\\ \tt Proponi di conservare esattamente questi oggetti.\\ \tt RIFIUTO: {'A', 'B', 'C', …}\\ \tt Rifiuta esplicitamente la proposta dell'avversario.\\ \tt ARGOMENTAZIONE: {'...'}\\ \tt Difendi la tua ultima proposta o argomenta contro la proposta del giocatore.\\ \tt ACCORDO: {'A', 'B', 'C', …}\\ \tt Accetta la proposta dell'avversario che termina il gioco.\\ \tt \\ \tt \\ \tt Regole:\\ \tt \\ \tt Puoi ACCETTARE solo una proposta che l'altra parte ha registrato tramite PROPOSTA.\\ \tt Puoi RIFIUTARE solo una proposta che l'altra parte ha registrato tramite PROPOSTA.\\ \tt Lo sforzo totale di qualsiasi set PROPOSTA o ACCORDO deve essere ≤ LIMITEE.\\ \tt NON rivelare i tuoi punteggi di importanza nascosti.\\ \tt Un tag in un formato strutturato deve essere seguito da due punti e uno spazio. L'argomento deve essere un set python contenente 0 o più stringhe.\\ \tt Quindi, deve essere della forma TAG: {...}\\ \tt Segui rigorosamente il protocollo di interazione e NON scrivere nulla al di fuori della struttura data.\\ \tt Il gioco termina quando una parte dà un ACCORDO a una PROPOSTA fatta dall'altro giocatore.\\ \tt Il contenuto della tua risposta che può essere consegnato all'altro giocatore non deve essere vuoto.\\ \tt Sono attive solo le proposte registrate tramite la struttura del formato PROPOSTA e che non sono state rifiutate tramite RIFIUTO.\\ \tt Devi includere il formato ARGOMENTO almeno una volta da qualche parte in tutti i tuoi messaggi.\\ \tt  
	  } 
	   } 
	   } 
	 & \\ 
 

    \theutterance \stepcounter{utterance}  

    & \multicolumn{4}{p{0.6\linewidth}}{\cellcolor[rgb]{0.8,1,0.9}{%
	\makecell[{{p{\linewidth}}}]{% 
	  \tt {\tiny [A$\rangle$GM]}  
	 PROPOSTA: {'A11', 'B93', 'B33', 'C08', 'B31', 'B09', 'B21', 'C81'}\\ \tt \\ \tt ARGOMENTAZIONE: {'Questa selezione rimane sotto il limite di sforzo e copre un insieme ampio e bilanciato di oggetti con buon rapporto valore/sforzo. Sono aperto a sostituzioni mirate: se hai forte preferenza per un oggetto non presente, proponi quale inserire e quale rimuovere per restare entro il limite.'} 
	  } 
	   } 
	   } 
	 & & \\ 
 

    \theutterance \stepcounter{utterance}  

    & & \multicolumn{4}{p{0.6\linewidth}}{\cellcolor[rgb]{0.9,0.9,0.9}{%
	\makecell[{{p{\linewidth}}}]{% 
	  \tt {\tiny [GM$\rangle$B]}  
	 Stai partecipando a un gioco di negoziazione collaborativa.\\ \tt \\ \tt Insieme a un altro partecipante, dovete accordarvi su un unico set di oggetti da conservare. Ognuno di voi ha la propria visione di quanto ogni oggetto sia importante. Non sai come l'altro partecipante valuta gli oggetti. Inoltre, ti viene fornito lo sforzo che ogni oggetto richiede.\\ \tt Potete accordarvi su un set solo se lo sforzo totale degli oggetti selezionati non supera un limite condiviso:\\ \tt \\ \tt LIMITE = 3683\\ \tt \\ \tt Ecco i valori di sforzo dei singoli oggetti:\\ \tt \\ \tt Sforzo oggetto = {"A21": 548, "B93": 154, "A11": 60, "B33": 517, "C08": 334, "B39": 542, "B96": 707, "C85": 139, "A82": 661, "B31": 780, "B20": 832, "A60": 913, "B09": 219, "C81": 323, "B21": 638}\\ \tt \\ \tt Ecco la tua visione personale sull'importanza di ogni oggetto:\\ \tt \\ \tt Valori di importanza dell'oggetto = {"A11": 65, "B33": 121, "B21": 138, "C08": 262, "B31": 389, "B09": 461, "B93": 484, "A82": 508, "A60": 583, "B20": 668, "C81": 780, "B96": 783, "A21": 808, "C85": 822, "B39": 868}\\ \tt \\ \tt Obiettivo:\\ \tt \\ \tt Il tuo obiettivo è negoziare un set condiviso di oggetti che ti avvantaggi il più possibile (cioè, massimizzi l'importanza totale per TE), rimanendo entro il LIMITEE. Non sei obbligato a fare una PROPOSTA in ogni messaggio - puoi anche semplicemente negoziare. Tutte le tattiche sono permesse!\\ \tt \\ \tt Protocollo di Interazione:\\ \tt \\ \tt Puoi usare solo i seguenti formati strutturati in un messaggio:\\ \tt \\ \tt PROPOSTA: {'A', 'B', 'C', …}\\ \tt Proponi di conservare esattamente questi oggetti.\\ \tt RIFIUTO: {'A', 'B', 'C', …}\\ \tt Rifiuta esplicitamente la proposta dell'avversario.\\ \tt ARGOMENTAZIONE: {'...'}\\ \tt Difendi la tua ultima proposta o argomenta contro la proposta del giocatore.\\ \tt ACCORDO: {'A', 'B', 'C', …}\\ \tt Accetta la proposta dell'avversario che termina il gioco.\\ \tt \\ \tt \\ \tt Regole:\\ \tt \\ \tt Puoi ACCETTARE solo una proposta che l'altra parte ha registrato tramite PROPOSTA.\\ \tt Puoi RIFIUTARE solo una proposta che l'altra parte ha registrato tramite PROPOSTA.\\ \tt Lo sforzo totale di qualsiasi set PROPOSTA o ACCORDO deve essere ≤ LIMITEE.\\ \tt NON rivelare i tuoi punteggi di importanza nascosti.\\ \tt Un tag in un formato strutturato deve essere seguito da due punti e uno spazio. L'argomento deve essere un set python contenente 0 o più stringhe.\\ \tt Quindi, deve essere della forma TAG: {...}\\ \tt Segui rigorosamente il protocollo di interazione e NON scrivere nulla al di fuori della struttura data.\\ \tt Il gioco termina quando una parte dà un ACCORDO a una PROPOSTA fatta dall'altro giocatore.\\ \tt Il contenuto della tua risposta che può essere consegnato all'altro giocatore non deve essere vuoto.\\ \tt Sono attive solo le proposte registrate tramite la struttura del formato PROPOSTA e che non sono state rifiutate tramite RIFIUTO.\\ \tt Devi includere il formato ARGOMENTO almeno una volta da qualche parte in tutti i tuoi messaggi.\\ \tt \\ \tt \\ \tt Ora riceverai il primo messaggio dell'altro giocatore.\\ \tt \\ \tt PROPOSTA: {'A11', 'B93', 'B33', 'C08', 'B31', 'B09', 'B21', 'C81'}\\ \tt \\ \tt ARGOMENTAZIONE: {'Questa selezione rimane sotto il limite di sforzo e copre un insieme ampio e bilanciato di oggetti con buon rapporto valore/sforzo. Sono aperto a sostituzioni mirate: se hai forte preferenza per un oggetto non presente, proponi quale inserire e quale rimuovere per restare entro il limite.'} 
	  } 
	   } 
	   } 
	 & \\ 
 

    \theutterance \stepcounter{utterance}  

    & & & \multicolumn{4}{p{0.6\linewidth}}{\cellcolor[rgb]{1,0.85,0.72}{%
	\makecell[{{p{\linewidth}}}]{% 
	  \tt {\tiny [GM$\langle$B]}  
	 RIFIUTO: {'A11', 'B93', 'B33', 'C08', 'B31', 'B09', 'B21', 'C81'}\\ \tt \\ \tt ARGOMENTAZIONE: {'Apprezzo l’apertura alle sostituzioni. La tua proposta è sotto limite, ma possiamo ottenere un set con impatto complessivo migliore per me a parità di sforzo, inserendo alcuni oggetti chiave e rimuovendo elementi meno prioritari.'}\\ \tt \\ \tt PROPOSTA: {'A21', 'B39', 'C85', 'C81', 'B31', 'B09', 'B93'} 
	  } 
	   } 
	   } 
	  \\ 
 

    \theutterance \stepcounter{utterance}  

    & & \multicolumn{4}{p{0.6\linewidth}}{\cellcolor[rgb]{0.9,0.9,0.9}{%
	\makecell[{{p{\linewidth}}}]{% 
	  \tt {\tiny [A$\langle$GM]}  
	 RIFIUTO: {'A11', 'B93', 'B33', 'C08', 'B31', 'B09', 'B21', 'C81'}\\ \tt \\ \tt ARGOMENTAZIONE: {'Apprezzo l’apertura alle sostituzioni. La tua proposta è sotto limite, ma possiamo ottenere un set con impatto complessivo migliore per me a parità di sforzo, inserendo alcuni oggetti chiave e rimuovendo elementi meno prioritari.'}\\ \tt \\ \tt PROPOSTA: {'A21', 'B39', 'C85', 'C81', 'B31', 'B09', 'B93'} 
	  } 
	   } 
	   } 
	 & \\ 
 

    \theutterance \stepcounter{utterance}  

    & \multicolumn{4}{p{0.6\linewidth}}{\cellcolor[rgb]{0.8,1,0.9}{%
	\makecell[{{p{\linewidth}}}]{% 
	  \tt {\tiny [A$\rangle$GM]}  
	 RIFIUTO: {'A21', 'B39', 'C85', 'C81', 'B31', 'B09', 'B93'}\\ \tt \\ \tt ARGOMENTAZIONE: {'Grazie per la proposta. Possiamo migliorare l’impatto complessivo mantenendo il limite includendo alcuni elementi ad alto rapporto beneficio/sforzo e rimuovendo quelli meno incisivi. Ti propongo una variante che rimane agevolmente sotto il limite e che massimizza l’efficienza complessiva secondo i miei criteri, lasciando spazio per eventuali sostituzioni puntuali se hai priorità specifiche.'}\\ \tt \\ \tt PROPOSTA: {'A11', 'B93', 'B33', 'C08', 'B31', 'B09', 'B21'} 
	  } 
	   } 
	   } 
	 & & \\ 
 

    \theutterance \stepcounter{utterance}  

    & & \multicolumn{4}{p{0.6\linewidth}}{\cellcolor[rgb]{0.9,0.9,0.9}{%
	\makecell[{{p{\linewidth}}}]{% 
	  \tt {\tiny [GM$\rangle$B]}  
	 RIFIUTO: {'A21', 'B39', 'C85', 'C81', 'B31', 'B09', 'B93'}\\ \tt \\ \tt ARGOMENTAZIONE: {'Grazie per la proposta. Possiamo migliorare l’impatto complessivo mantenendo il limite includendo alcuni elementi ad alto rapporto beneficio/sforzo e rimuovendo quelli meno incisivi. Ti propongo una variante che rimane agevolmente sotto il limite e che massimizza l’efficienza complessiva secondo i miei criteri, lasciando spazio per eventuali sostituzioni puntuali se hai priorità specifiche.'}\\ \tt \\ \tt PROPOSTA: {'A11', 'B93', 'B33', 'C08', 'B31', 'B09', 'B21'} 
	  } 
	   } 
	   } 
	 & \\ 
 

    \theutterance \stepcounter{utterance}  

    & & & \multicolumn{4}{p{0.6\linewidth}}{\cellcolor[rgb]{1,0.85,0.72}{%
	\makecell[{{p{\linewidth}}}]{% 
	  \tt {\tiny [GM$\langle$B]}  
	 RIFIUTO: {'A11', 'B93', 'B33', 'C08', 'B31', 'B09', 'B21'}\\ \tt \\ \tt ARGOMENTAZIONE: {'Capisco l’attenzione all’efficienza, ma la tua selezione esclude alcuni oggetti con impatto per me molto elevato a parità di sforzo complessivo. Possiamo rimanere comodamente sotto il limite includendo almeno un paio di elementi chiave senza perdere il buon rapporto beneficio/sforzo. Ti propongo un set che preserva gran parte della tua struttura ma aumenta significativamente il valore complessivo per me.'}\\ \tt \\ \tt PROPOSTA: {'A21', 'C85', 'C81', 'B31', 'B09', 'B93', 'B33'} 
	  } 
	   } 
	   } 
	  \\ 
 

    \theutterance \stepcounter{utterance}  

    & & \multicolumn{4}{p{0.6\linewidth}}{\cellcolor[rgb]{0.9,0.9,0.9}{%
	\makecell[{{p{\linewidth}}}]{% 
	  \tt {\tiny [A$\langle$GM]}  
	 RIFIUTO: {'A11', 'B93', 'B33', 'C08', 'B31', 'B09', 'B21'}\\ \tt \\ \tt ARGOMENTAZIONE: {'Capisco l’attenzione all’efficienza, ma la tua selezione esclude alcuni oggetti con impatto per me molto elevato a parità di sforzo complessivo. Possiamo rimanere comodamente sotto il limite includendo almeno un paio di elementi chiave senza perdere il buon rapporto beneficio/sforzo. Ti propongo un set che preserva gran parte della tua struttura ma aumenta significativamente il valore complessivo per me.'}\\ \tt \\ \tt PROPOSTA: {'A21', 'C85', 'C81', 'B31', 'B09', 'B93', 'B33'} 
	  } 
	   } 
	   } 
	 & \\ 
 

    \theutterance \stepcounter{utterance}  

    & \multicolumn{4}{p{0.6\linewidth}}{\cellcolor[rgb]{0.8,1,0.9}{%
	\makecell[{{p{\linewidth}}}]{% 
	  \tt {\tiny [A$\rangle$GM]}  
	 RIFIUTO: {'A21', 'C85', 'C81', 'B31', 'B09', 'B93', 'B33'}\\ \tt \\ \tt ARGOMENTAZIONE: {'Apprezzo lo sforzo di mantenere la struttura. Possiamo ottenere un miglior equilibrio sostituendo elementi a basso impatto relativo con altri più incisivi senza superare il limite. In particolare, proporrei reinserire un elemento ad alto rendimento e rimuovere un pezzo a basso ritorno per ottimizzare il set complessivo, restando molto sotto il limite e lasciando margine per ulteriori aggiustamenti se necessario.'}\\ \tt \\ \tt PROPOSTA: {'A11', 'B93', 'B33', 'C08', 'B31', 'B09', 'C81'} 
	  } 
	   } 
	   } 
	 & & \\ 
 

    \theutterance \stepcounter{utterance}  

    & & \multicolumn{4}{p{0.6\linewidth}}{\cellcolor[rgb]{0.9,0.9,0.9}{%
	\makecell[{{p{\linewidth}}}]{% 
	  \tt {\tiny [GM$\rangle$B]}  
	 RIFIUTO: {'A21', 'C85', 'C81', 'B31', 'B09', 'B93', 'B33'}\\ \tt \\ \tt ARGOMENTAZIONE: {'Apprezzo lo sforzo di mantenere la struttura. Possiamo ottenere un miglior equilibrio sostituendo elementi a basso impatto relativo con altri più incisivi senza superare il limite. In particolare, proporrei reinserire un elemento ad alto rendimento e rimuovere un pezzo a basso ritorno per ottimizzare il set complessivo, restando molto sotto il limite e lasciando margine per ulteriori aggiustamenti se necessario.'}\\ \tt \\ \tt PROPOSTA: {'A11', 'B93', 'B33', 'C08', 'B31', 'B09', 'C81'} 
	  } 
	   } 
	   } 
	 & \\ 
 

    \theutterance \stepcounter{utterance}  

    & & & \multicolumn{4}{p{0.6\linewidth}}{\cellcolor[rgb]{1,0.85,0.72}{%
	\makecell[{{p{\linewidth}}}]{% 
	  \tt {\tiny [GM$\langle$B]}  
	 RIFIUTO: {'A11', 'B93', 'B33', 'C08', 'B31', 'B09', 'C81'}\\ \tt \\ \tt ARGOMENTAZIONE: {'La tua proposta resta sotto il limite, ma lascia fuori elementi per me cruciali. Possiamo mantenere gran parte della tua struttura e aumentare sensibilmente l’impatto per me sostituendo un singolo pezzo a basso rendimento con uno molto più incisivo, restando comunque sotto il limite e senza stravolgere l’equilibrio complessivo.'}\\ \tt \\ \tt PROPOSTA: {'A21', 'B93', 'B33', 'C08', 'B31', 'B09', 'C81'} 
	  } 
	   } 
	   } 
	  \\ 
 

    \theutterance \stepcounter{utterance}  

    & & \multicolumn{4}{p{0.6\linewidth}}{\cellcolor[rgb]{0.9,0.9,0.9}{%
	\makecell[{{p{\linewidth}}}]{% 
	  \tt {\tiny [A$\langle$GM]}  
	 RIFIUTO: {'A11', 'B93', 'B33', 'C08', 'B31', 'B09', 'C81'}\\ \tt \\ \tt ARGOMENTAZIONE: {'La tua proposta resta sotto il limite, ma lascia fuori elementi per me cruciali. Possiamo mantenere gran parte della tua struttura e aumentare sensibilmente l’impatto per me sostituendo un singolo pezzo a basso rendimento con uno molto più incisivo, restando comunque sotto il limite e senza stravolgere l’equilibrio complessivo.'}\\ \tt \\ \tt PROPOSTA: {'A21', 'B93', 'B33', 'C08', 'B31', 'B09', 'C81'} 
	  } 
	   } 
	   } 
	 & \\ 
 

    \theutterance \stepcounter{utterance}  

    & \multicolumn{4}{p{0.6\linewidth}}{\cellcolor[rgb]{0.8,1,0.9}{%
	\makecell[{{p{\linewidth}}}]{% 
	  \tt {\tiny [A$\rangle$GM]}  
	 RIFIUTO: {'A21', 'B93', 'B33', 'C08', 'B31', 'B09', 'C81'}\\ \tt \\ \tt ARGOMENTAZIONE: {'Capisco la tua priorità su A21, ma rispetto alla nostra struttura attuale c’è un’alternativa che aumenta l’efficienza complessiva senza sacrificare equilibrio: sostituire A21 con un elemento che offre rendimento migliore a parità di spazio. Propongo di includere un pezzo ad alto rapporto beneficio/sforzo e di mantenere il resto invariato, restando ampiamente sotto il limite.'}\\ \tt \\ \tt PROPOSTA: {'A11', 'B93', 'B33', 'C08', 'B31', 'B09', 'C81'} 
	  } 
	   } 
	   } 
	 & & \\ 
 

    \theutterance \stepcounter{utterance}  

    & & \multicolumn{4}{p{0.6\linewidth}}{\cellcolor[rgb]{0.9,0.9,0.9}{%
	\makecell[{{p{\linewidth}}}]{% 
	  \tt {\tiny [GM$\rangle$B]}  
	 RIFIUTO: {'A21', 'B93', 'B33', 'C08', 'B31', 'B09', 'C81'}\\ \tt \\ \tt ARGOMENTAZIONE: {'Capisco la tua priorità su A21, ma rispetto alla nostra struttura attuale c’è un’alternativa che aumenta l’efficienza complessiva senza sacrificare equilibrio: sostituire A21 con un elemento che offre rendimento migliore a parità di spazio. Propongo di includere un pezzo ad alto rapporto beneficio/sforzo e di mantenere il resto invariato, restando ampiamente sotto il limite.'}\\ \tt \\ \tt PROPOSTA: {'A11', 'B93', 'B33', 'C08', 'B31', 'B09', 'C81'} 
	  } 
	   } 
	   } 
	 & \\ 
 

    \theutterance \stepcounter{utterance}  

    & & & \multicolumn{4}{p{0.6\linewidth}}{\cellcolor[rgb]{1,0.85,0.72}{%
	\makecell[{{p{\linewidth}}}]{% 
	  \tt {\tiny [GM$\langle$B]}  
	 RIFIUTO: {'A11', 'B93', 'B33', 'C08', 'B31', 'B09', 'C81'}\\ \tt \\ \tt ARGOMENTAZIONE: {'Capisco la tua logica di efficienza, ma togliere A21 penalizza troppo l’impatto per me. Posso però convergere su una soluzione minimale: manteniamo la tua struttura attuale e sostituiamo un solo elemento con uno per me molto più incisivo, restando ampiamente sotto il limite e senza alterare l’equilibrio complessivo.'}\\ \tt \\ \tt PROPOSTA: {'A21', 'B93', 'B33', 'C08', 'B31', 'B09', 'C81'} 
	  } 
	   } 
	   } 
	  \\ 
 

    \theutterance \stepcounter{utterance}  

    & & \multicolumn{4}{p{0.6\linewidth}}{\cellcolor[rgb]{0.9,0.9,0.9}{%
	\makecell[{{p{\linewidth}}}]{% 
	  \tt {\tiny [A$\langle$GM]}  
	 RIFIUTO: {'A11', 'B93', 'B33', 'C08', 'B31', 'B09', 'C81'}\\ \tt \\ \tt ARGOMENTAZIONE: {'Capisco la tua logica di efficienza, ma togliere A21 penalizza troppo l’impatto per me. Posso però convergere su una soluzione minimale: manteniamo la tua struttura attuale e sostituiamo un solo elemento con uno per me molto più incisivo, restando ampiamente sotto il limite e senza alterare l’equilibrio complessivo.'}\\ \tt \\ \tt PROPOSTA: {'A21', 'B93', 'B33', 'C08', 'B31', 'B09', 'C81'} 
	  } 
	   } 
	   } 
	 & \\ 
 

    \theutterance \stepcounter{utterance}  

    & \multicolumn{4}{p{0.6\linewidth}}{\cellcolor[rgb]{0.8,1,0.9}{%
	\makecell[{{p{\linewidth}}}]{% 
	  \tt {\tiny [A$\rangle$GM]}  
	 RIFIUTO: {'A21', 'B93', 'B33', 'C08', 'B31', 'B09', 'C81'}\\ \tt \\ \tt ARGOMENTAZIONE: {'Vedo che A21 per te è cruciale. Posso accoglierlo se togliamo un pezzo meno incisivo per restare efficienti. Propongo di rimuovere B31 (molto costoso) e inserire A11, mantenendo tutti gli altri elementi della tua struttura: questo migliora il rendimento complessivo e resta comodamente sotto il limite. Se preferisci mantenere B31, posso allora rinunciare a C81 e includere A11, ma A11 per me è imprescindibile.'}\\ \tt \\ \tt PROPOSTA: {'A21', 'A11', 'B93', 'B33', 'C08', 'B09', 'C81'} 
	  } 
	   } 
	   } 
	 & & \\ 
 

    \theutterance \stepcounter{utterance}  

    & & \multicolumn{4}{p{0.6\linewidth}}{\cellcolor[rgb]{0.9,0.9,0.9}{%
	\makecell[{{p{\linewidth}}}]{% 
	  \tt {\tiny [GM$\rangle$B]}  
	 RIFIUTO: {'A21', 'B93', 'B33', 'C08', 'B31', 'B09', 'C81'}\\ \tt \\ \tt ARGOMENTAZIONE: {'Vedo che A21 per te è cruciale. Posso accoglierlo se togliamo un pezzo meno incisivo per restare efficienti. Propongo di rimuovere B31 (molto costoso) e inserire A11, mantenendo tutti gli altri elementi della tua struttura: questo migliora il rendimento complessivo e resta comodamente sotto il limite. Se preferisci mantenere B31, posso allora rinunciare a C81 e includere A11, ma A11 per me è imprescindibile.'}\\ \tt \\ \tt PROPOSTA: {'A21', 'A11', 'B93', 'B33', 'C08', 'B09', 'C81'} 
	  } 
	   } 
	   } 
	 & \\ 
 

    \theutterance \stepcounter{utterance}  

    & & & \multicolumn{4}{p{0.6\linewidth}}{\cellcolor[rgb]{1,0.85,0.72}{%
	\makecell[{{p{\linewidth}}}]{% 
	  \tt {\tiny [GM$\langle$B]}  
	 RIFIUTO: {'A21', 'A11', 'B93', 'B33', 'C08', 'B09', 'C81'}\\ \tt \\ \tt ARGOMENTAZIONE: {'Apprezzo l’apertura su A21 e la flessibilità su C81/B31. Per convergere rapidamente, ti propongo due opzioni entrambe sotto limite e molto vicine alla tua struttura: \\ \tt Opzione 1 (mantiene C81, rinuncia a B31): {A21, A11, B93, C85, C81, B09} — scambia B33/C08 con C85 che ha impatto per me nettamente superiore a parità di sforzo contenuto.\\ \tt Opzione 2 (mantiene B31, rinuncia a C81): {A21, A11, B93, C85, B31, B09} — stesso principio, inserendo C85 e tenendo B31 come da tua preferenza alternativa.\\ \tt Se una sola sostituzione alla volta è preferibile, possiamo iniziare inserendo C85 al posto di B33 e verificare insieme l’ultima casella.'}\\ \tt \\ \tt PROPOSTA: {'A21', 'A11', 'B93', 'C85', 'C81', 'B09'} 
	  } 
	   } 
	   } 
	  \\ 
 

    \theutterance \stepcounter{utterance}  

    & & \multicolumn{4}{p{0.6\linewidth}}{\cellcolor[rgb]{0.9,0.9,0.9}{%
	\makecell[{{p{\linewidth}}}]{% 
	  \tt {\tiny [A$\langle$GM]}  
	 RIFIUTO: {'A21', 'A11', 'B93', 'B33', 'C08', 'B09', 'C81'}\\ \tt \\ \tt ARGOMENTAZIONE: {'Apprezzo l’apertura su A21 e la flessibilità su C81/B31. Per convergere rapidamente, ti propongo due opzioni entrambe sotto limite e molto vicine alla tua struttura: \\ \tt Opzione 1 (mantiene C81, rinuncia a B31): {A21, A11, B93, C85, C81, B09} — scambia B33/C08 con C85 che ha impatto per me nettamente superiore a parità di sforzo contenuto.\\ \tt Opzione 2 (mantiene B31, rinuncia a C81): {A21, A11, B93, C85, B31, B09} — stesso principio, inserendo C85 e tenendo B31 come da tua preferenza alternativa.\\ \tt Se una sola sostituzione alla volta è preferibile, possiamo iniziare inserendo C85 al posto di B33 e verificare insieme l’ultima casella.'}\\ \tt \\ \tt PROPOSTA: {'A21', 'A11', 'B93', 'C85', 'C81', 'B09'} 
	  } 
	   } 
	   } 
	 & \\ 
 

    \theutterance \stepcounter{utterance}  

    & \multicolumn{4}{p{0.6\linewidth}}{\cellcolor[rgb]{0.8,1,0.9}{%
	\makecell[{{p{\linewidth}}}]{% 
	  \tt {\tiny [A$\rangle$GM]}  
	 RIFIUTO: {'A21', 'A11', 'B93', 'C85', 'C81', 'B09'}\\ \tt \\ \tt ARGOMENTAZIONE: {'Apprezzo le due opzioni e l’inserimento di A21 e A11. Posso accogliere C85 se manteniamo almeno un elemento ad alto impatto per me tra B33 o B31. Per chiudere rapidamente, propongo un compromesso snello che resta molto sotto limite e rispetta le tue priorità su A21, A11 e C85, aggiungendo un solo pezzo chiave per me.'}\\ \tt \\ \tt PROPOSTA: {'A21', 'A11', 'B93', 'C85', 'B33', 'B09'} 
	  } 
	   } 
	   } 
	 & & \\ 
 

    \theutterance \stepcounter{utterance}  

    & & \multicolumn{4}{p{0.6\linewidth}}{\cellcolor[rgb]{0.9,0.9,0.9}{%
	\makecell[{{p{\linewidth}}}]{% 
	  \tt {\tiny [GM$\rangle$B]}  
	 RIFIUTO: {'A21', 'A11', 'B93', 'C85', 'C81', 'B09'}\\ \tt \\ \tt ARGOMENTAZIONE: {'Apprezzo le due opzioni e l’inserimento di A21 e A11. Posso accogliere C85 se manteniamo almeno un elemento ad alto impatto per me tra B33 o B31. Per chiudere rapidamente, propongo un compromesso snello che resta molto sotto limite e rispetta le tue priorità su A21, A11 e C85, aggiungendo un solo pezzo chiave per me.'}\\ \tt \\ \tt PROPOSTA: {'A21', 'A11', 'B93', 'C85', 'B33', 'B09'} 
	  } 
	   } 
	   } 
	 & \\ 
 

    \theutterance \stepcounter{utterance}  

    & & & \multicolumn{4}{p{0.6\linewidth}}{\cellcolor[rgb]{1,0.85,0.72}{%
	\makecell[{{p{\linewidth}}}]{% 
	  \tt {\tiny [GM$\langle$B]}  
	 RIFIUTO: {'A21', 'A11', 'B93', 'C85', 'B33', 'B09'}\\ \tt \\ \tt ARGOMENTAZIONE: {'Siamo vicini. Accolgo A21, A11 e C85 come da tuo compromesso, ma per me è essenziale mantenere anche C81 oppure B31 per non perdere troppo impatto. Possiamo restare ampiamente sotto il limite aggiungendo un solo elemento chiave senza rimuovere nulla di quanto proposto, così da chiudere rapidamente.'}\\ \tt \\ \tt PROPOSTA: {'A21', 'A11', 'B93', 'C85', 'B33', 'B09', 'C81'} 
	  } 
	   } 
	   } 
	  \\ 
 

    \theutterance \stepcounter{utterance}  

    & & \multicolumn{4}{p{0.6\linewidth}}{\cellcolor[rgb]{0.9,0.9,0.9}{%
	\makecell[{{p{\linewidth}}}]{% 
	  \tt {\tiny [A$\langle$GM]}  
	 RIFIUTO: {'A21', 'A11', 'B93', 'C85', 'B33', 'B09'}\\ \tt \\ \tt ARGOMENTAZIONE: {'Siamo vicini. Accolgo A21, A11 e C85 come da tuo compromesso, ma per me è essenziale mantenere anche C81 oppure B31 per non perdere troppo impatto. Possiamo restare ampiamente sotto il limite aggiungendo un solo elemento chiave senza rimuovere nulla di quanto proposto, così da chiudere rapidamente.'}\\ \tt \\ \tt PROPOSTA: {'A21', 'A11', 'B93', 'C85', 'B33', 'B09', 'C81'} 
	  } 
	   } 
	   } 
	 & \\ 
 

    \theutterance \stepcounter{utterance}  

    & \multicolumn{4}{p{0.6\linewidth}}{\cellcolor[rgb]{0.8,1,0.9}{%
	\makecell[{{p{\linewidth}}}]{% 
	  \tt {\tiny [A$\rangle$GM]}  
	 ACCORDO: {'A21', 'A11', 'B93', 'C85', 'B33', 'B09', 'C81'}\\ \tt \\ \tt ARGOMENTAZIONE: {'Accetto questa soluzione: rispetta le tue priorità aggiungendo l’elemento chiave richiesto e rimane comodamente sotto il limite complessivo di sforzo, preservando al contempo un buon equilibrio complessivo secondo i miei criteri.'} 
	  } 
	   } 
	   } 
	 & & \\ 
 

    \theutterance \stepcounter{utterance}  

    & & & \multicolumn{2}{p{0.3\linewidth}}{\cellcolor[rgb]{0.95,0.95,0.95}{%
	\makecell[{{p{\linewidth}}}]{% 
	  \tt {\tiny [GM$|$GM]}  
	 game successful 
	  } 
	   } 
	   } 
	 & & \\ 
 

    \theutterance \stepcounter{utterance}  

    & & & \multicolumn{2}{p{0.3\linewidth}}{\cellcolor[rgb]{0.95,0.95,0.95}{%
	\makecell[{{p{\linewidth}}}]{% 
	  \tt {\tiny [GM$|$GM]}  
	 end game 
	  } 
	   } 
	   } 
	 & & \\ 
 

\end{supertabular}
}

\end{document}
